%% kapa: exam=1 round=36 year=2025 period=1 form=A subject=economics subject_ko=경제학원론 seq=36 number=76
%% source: Q-Net 감정평가사 기출문제 (cst00309) artlSeq=5231525 file=raw/1008/5231525/2025년 제36회 감정평가사 1차시험 1교시 문제지 원본.pdf
%% conversion: kordoc 4.24.51 -> tools/gen1.py (auto); review=proofread-sonnet; flags=none
\begin{question}{76}
소비이론 중 상대소득가설에 관한 설명으로 옳은 것은?
\begin{choices}
  \item 현재 소비는 다른 사람의 소비와 자신의 과거 소비에 의존한다.
  \item 유산이 증가하면 항상소비는 증가한다.
  \item 갑자기 받은 현금성보조금으로 항상소비는 증가한다.
  \item 소비에는 전시효과(demonstration effect)가 존재하지 않는다.
  \item 노인이 되면 소비는 노후소득에 따라 결정된다.
\end{choices}
\end{question}
\answer{1}
